\begin{blocksection}
Gleam is a functional programming language, a paradigm where programs are built
by applying and composing functions, and function definitions are trees of
expressions mapping values to other values rather than sequences of statements
that update program state.

Gleam is also statically-typed and type-safe. Type safety is the extent to which
a language prevents type errors, such as adding a string to an integer. Type
enforcement can be static (at compile time), dynamic (at run-time), or both. In
static type systems, type safety usually guarantees that an expression's
eventual value belongs to its static type.

When writing functions that operate on a given input, Gleam requires that you
handle any theoretically possible input case, even if you've potentially handled
such a case earlier. For instance, the function \texttt{list.max} (which behaves
very similarly to Python's \texttt{max} function) does not return an integer
directly, but a \texttt{Result}. This covers an empty list being passed into \texttt{max}: it returns
\texttt{Ok(value)} or \texttt{Error(Nil)} even if you've already ruled out the
empty case, guaranteeing an \texttt{Ok} result.

Similarly, a \texttt{case} statement in Gleam must handle every theoretically possible case, 
even if that case may not be possible in the context
of your own code. For instance, say you have a function that exclusively outputs
even integers, and you use one of those integers \texttt{n} in a \texttt{case}
statement:
\begin{lstlisting}
case n % 2 {
    0 -> //some code for the case in which n is even
    _ -> //some code for the case in which n is odd
}
\end{lstlisting}
Gleam will require you to include the catch-all \texttt{\_} case, regardless of
whether \texttt{n} could be odd. Even writing a \texttt{1} case is not enough:
Gleam cannot prove that \texttt{n \% 2} is only ever \texttt{0} or \texttt{1}
(it is \texttt{-1} for negative odd \texttt{n}), so an \texttt{Int} case always needs \texttt{\_}.
\newline
\end{blocksection}